% GENERATED FILE. Edit canonical lessons, then run npm run generate:books.
% canonical-source-hash: fnv1a64:30d0763a081b046a
% canonical-lessons: BN-C105-jhara, BN-C105-gochhano, BN-C105-sajano, BN-C105-bhaja, BN-C105-meshano

\chapter{To dust, To tidy up, To decorate, To fry, To mix}
\label{ch:bn-to-dust-to-tidy-up-to-decorate-to-fry-to-mix}
\hlchaptermodality{105}

\begin{chapteropening}
\textbf{By the end of this chapter:} \emph{I can say to dust, to tidy up, to decorate, to fry and to mix.}

Complete the last lesson of chapter 105: I can say to dust, to tidy up, to decorate, to fry and to mix.
\end{chapteropening}

\section[\texorpdfstring{\bn{ঝাড়া} (jhāṛā)}{ঝাড়া (jhāṛā)}]{\bn{ঝাড়া} (jhāṛā) — to dust, to shake out}

\label{lesson:BN-C105-jhara}



\begin{quote}
Before the new one: say the Bengali for closed, then the Bengali for to wipe.
\end{quote}

\subsection*{You'll want to know: \bn{ঝাড়া}}
\textbf{\bn{ঝাড়া}} — \emph{jhāṛā} — \textquotedblleft{}to dust, to shake out\textquotedblright{}.

\begin{grammarlens}[title={What you've built}]
One more word for this chapter.
\end{grammarlens}

\subsection*{Your turn}
\begin{itemize}
\raggedright
  \item \emph{Say it:} \emph{jhāṛā}
  \item \emph{Say it:} \emph{jhāṛā}, once more
  \item \emph{Say it:} say \emph{jhāṛā}
  \item \emph{Recall:} say the Bengali for closed, then the Bengali for to wipe, then say \emph{jhāṛā} again
\end{itemize}

\begin{tcolorbox}[breakable,colback=teal!4,colframe=teal!35!black,title={Before you move on}]
What does \bn{ঝাড়া} mean? (\textquotedblleft{}To dust, to shake out\textquotedblright{}.) Say it once more.
\end{tcolorbox}
\section[\texorpdfstring{\bn{গোছানো} (gochhāno)}{গোছানো (gochhāno)}]{\bn{গোছানো} (gochhāno) — to tidy up}

\label{lesson:BN-C105-gochhano}



\begin{quote}
Before the new one: say the Bengali for to wipe, then the Bengali for to dust.
\end{quote}

\subsection*{You'll want to know: \bn{গোছানো}}
\textbf{\bn{গোছানো}} — \emph{gochhāno} — \textquotedblleft{}to tidy up\textquotedblright{}.

\begin{grammarlens}[title={What you've built}]
One more word for this chapter.
\end{grammarlens}

\subsection*{Your turn}
\begin{itemize}
\raggedright
  \item \emph{Say it:} \emph{gochhāno}
  \item \emph{Say it:} \emph{gochhāno}, once more
  \item \emph{Say it:} say \emph{gochhāno}
  \item \emph{Recall:} say the Bengali for to wipe, then the Bengali for to dust, then say \emph{gochhāno} again
\end{itemize}

\begin{tcolorbox}[breakable,colback=teal!4,colframe=teal!35!black,title={Before you move on}]
What does \bn{গোছানো} mean? (\textquotedblleft{}To tidy up\textquotedblright{}.) Say it once more.
\end{tcolorbox}
\section[\texorpdfstring{\bn{সাজানো} (sājāno)}{সাজানো (sājāno)}]{\bn{সাজানো} (sājāno) — to decorate, to arrange}

\label{lesson:BN-C105-sajano}



\begin{quote}
Before the new one: say the Bengali for to dust, then the Bengali for to tidy up.
\end{quote}

\subsection*{You'll want to know: \bn{সাজানো}}
\textbf{\bn{সাজানো}} — \emph{sājāno} — \textquotedblleft{}to decorate, to arrange\textquotedblright{}.

\begin{grammarlens}[title={What you've built}]
One more word for this chapter.
\end{grammarlens}

\subsection*{Your turn}
\begin{itemize}
\raggedright
  \item \emph{Say it:} \emph{sājāno}
  \item \emph{Say it:} \emph{sājāno}, once more
  \item \emph{Say it:} say \emph{sājāno}
  \item \emph{Recall:} say the Bengali for to dust, then the Bengali for to tidy up, then say \emph{sājāno} again
\end{itemize}

\begin{tcolorbox}[breakable,colback=teal!4,colframe=teal!35!black,title={Before you move on}]
What does \bn{সাজানো} mean? (\textquotedblleft{}To decorate, to arrange\textquotedblright{}.) Say it once more.
\end{tcolorbox}
\section[\texorpdfstring{\bn{ভাজা} (bhājā)}{ভাজা (bhājā)}]{\bn{ভাজা} (bhājā) — to fry}

\label{lesson:BN-C105-bhaja}



\begin{quote}
Before the new one: say the Bengali for to tidy up, then the Bengali for to decorate.

Which two Bengali digits lie? (\textbf{\bn{৪}} looks like an 8 and means four; \textbf{\bn{৭}} looks like a 9 and means seven.)
\end{quote}

\subsection*{You'll want to know: \bn{ভাজা}}
\textbf{\bn{ভাজা}} — \emph{bhājā} — \textquotedblleft{}to fry\textquotedblright{}.

\begin{grammarlens}[title={What you've built}]
One more word for this chapter.
\end{grammarlens}

\subsection*{Your turn}
\begin{itemize}
\raggedright
  \item \emph{Say it:} \emph{bhājā}
  \item \emph{Say it:} \emph{bhājā}, once more
  \item \emph{Say it:} say \emph{bhājā}
  \item \emph{Recall:} say the Bengali for to tidy up, then the Bengali for to decorate, then say \emph{bhājā} again
\end{itemize}

\begin{tcolorbox}[breakable,colback=teal!4,colframe=teal!35!black,title={Before you move on}]
What does \bn{ভাজা} mean? (\textquotedblleft{}To fry\textquotedblright{}.) Say it once more.
\end{tcolorbox}
\section[\texorpdfstring{\bn{মেশানো} (meshāno)}{মেশানো (meshāno)}]{\bn{মেশানো} (meshāno) — to mix}

\label{lesson:BN-C105-meshano}



\begin{quote}
Before the new one: say the Bengali for to decorate, then the Bengali for to fry.
\end{quote}

\subsection*{You'll want to know: \bn{মেশানো}}
\textbf{\bn{মেশানো}} — \emph{meshāno} — \textquotedblleft{}to mix\textquotedblright{}.

\begin{grammarlens}[title={What you've built}]
One more word for this chapter.
\end{grammarlens}

\subsection*{Your turn}
\begin{itemize}
\raggedright
  \item \emph{Say it:} \emph{meshāno}
  \item \emph{Say it:} \emph{meshāno}, once more
  \item \emph{Say it:} say \emph{meshāno}
  \item \emph{Recall:} say the Bengali for to decorate, then the Bengali for to fry, then say \emph{meshāno} again
\end{itemize}

\begin{tcolorbox}[breakable,colback=teal!4,colframe=teal!35!black,title={Before you move on}]
What does \bn{মেশানো} mean? (\textquotedblleft{}To mix\textquotedblright{}.) Say it once more.
\end{tcolorbox}
